\title{xLSTM: Extended Long Short-Term Memory}

\begin{abstract}
The foundational concepts of the constant error carousel and gating, introduced in the 1990s, have long underpinned the Long Short-Term Memory (LSTM) architecture. LSTMs have played a significant role in the advancement of deep learning, notably as the core component in the earliest Large Language Models (LLMs). However, with the emergence of Transformer architectures—characterized by their scalable, parallelizable self-attention mechanisms—a paradigm shift occurred, relegating LSTMs from their leading position. This paper asks: What can be achieved in language modeling by scaling LSTMs up to the billion-parameter regime, incorporating advanced methodologies developed for contemporary LLMs, while also addressing established limitations of LSTMs? To this end, we first introduce exponential gating, supplemented by normalization and stabilization strategies. We then rework the LSTM memory design and propose two variants: (i) sLSTM, featuring scalar-valued memory and updates with a revised memory mixing approach, and (ii) mLSTM, which offers a fully parallelizable structure via matrix memory and utilizes a covariance-based update mechanism. Embedding these extensions into the architecture of residual blocks results in xLSTM modules, which are composed residually to build up xLSTM networks. The use of exponential gating in combination with these reconfigured memory systems significantly enhances the capabilities of xLSTM models, enabling them to perform competitively with state-of-the-art Transformers and State Space Models in terms of both accuracy and scalability.\\
Code available at: \url{https://github.com/NX-AI/xlstm}
\end{abstract}

\section{Introduction}
\label{sec:intro}

The Long Short-Term Memory (LSTM) ideas~\citep{Hochreiter:91,Hochreiter:97nips,Hochreiter:97}, i.e., the constant error carousel and gating, were introduced to overcome the vanishing gradient problem of recurrent neural networks~\citep{Hochreiter:91,Hochreiter:00book}:

\begin{align}
\label{eq:lstm_idea}
  c_t \ = \ f_t \ c_{t-1} \ + \ 
          i_t \  z_t \ , \quad  
          h_t \ = \ o_t \ \psi({c_t}) \ . 
\end{align}

The constant error carousel operates by incrementally adjusting the cell state $c_{t-1}$ (green) through the addition of cell inputs $z_t$, with these updates regulated by sigmoid gates (blue). The process is managed by the input gate $i_t$ and forget gate $f_t$, which determine how the cell state is updated, whereas the output gate $o_t$ governs what portion of the cell's memory is revealed as the hidden state $h_t$. Prior to generating the hidden state, the cell state undergoes normalization or compression via $\psi$, after which the output gate mediates the hidden state’s production.

LSTMs have demonstrated their utility across a wide array of fields \citep{Hochreiter:01,Hochreiter:07,Schmidhuber:15}, dominating text generation tasks until the introduction of Transformers in 2017~\citep{Vaswani:17}. Their capability has been proven in many sequence-processing problems, including text generation~\citep{Graves:13,Karpathy:15blog}, handwriting synthesis~\citep{Graves:13}, sequence-to-sequence translation~\citep{Sutskever:14nips}, program evaluation~\citep{Zaremba:14arxiv}, image captioning~\citep{Karpathy:15,Hossain:19}, code generation~\citep{Karpathy:15blog}, hydrological modeling for rainfall-runoff~\citep{Kratzert:18,Kratzert:19}, and flood prediction models~\citep{Nearing:24}. Within reinforcement learning, LSTMs are regarded as the top-performing sequence models, having been utilized in state-of-the-art systems such as the AlphaStar agent for StarCraft II~\citep{Vinyals:17short}, OpenAI Five for Dota 2~\citep{OpenAI:19}, and magnetic controllers designed for nuclear fusion~\citep{Degrave:22short}. Their strength lies in their ability to form high-level abstractions, skillfully capturing semantic meaning within their cell states~\citep{Karpathy:15blog}, as highlighted by the discovery of number and syntax neurons~\citep{Lakretz:19}, linguistic neurons~\citep{Bau:19}, and sentiment neurons~\citep{Radford:17arxiv}. LSTMs continue to play a pivotal role in important contemporary applications~\citep{Degrave:22short,Nearing:24} and have maintained their relevance over time.

Although LSTMs have achieved notable accomplishments, they present three principal drawbacks: (i) Inability to update previously stored information. This is demonstrated with the {\it Nearest Neighbor Search} task (see also Appendix~\ref{sec:appNearestNeighborSearch}): Given a reference vector, the sequence must be examined in order to identify the most similar vector and output its corresponding value at the end of the sequence. As depicted in the left panel of Figure~\ref{fig:lstmProblems}, LSTMs show high mean squared error in this context, as they are unable to replace a previously retained value when a better match appears later in the sequence. Our proposed xLSTM overcomes this via exponential gating. (ii) Restricted memory capacity, necessitating that all information be encoded into scalar-valued cell states. This constraint is illustrated through the {\it Rare Token Prediction} scenario. In the right panel of Figure~\ref{fig:lstmProblems}, token prediction perplexity is reported on Wikitext-103~\citep{Merity:17} for groups based on token frequency. The LSTM demonstrates inferior performance for infrequent tokens, underscoring the impact of limited storage. The xLSTM addresses this issue using a memory organized as matrices. (iii) Inherent sequentiality due to hidden-to-hidden connections, commonly referred to as memory mixing, which makes parallel processing infeasible.

The aforementioned shortcomings have contributed to the dominance of the Transformer architecture~\citep{Vaswani:17} for language modeling applications. This raises an intriguing question: If these issues are resolved and LSTMs are scaled up to match the size of state-of-the-art Large Language Models, what levels of language modeling performance become possible?

\section{Extended Long Short-Term Memory}
\label{sec:xLSTM}

To address the shortcomings of LSTM, the Extended Long Short-Term Memory (xLSTM) framework implements two significant enhancements to the original LSTM concept delineated in Equation~\eqref{eq:lstm_idea}. These enhancements—namely, the use of exponential gating and the development of new memory architectures—yield two additional LSTM variants: (i) the sLSTM (refer to Section~\ref{sec:sLSTM}), characterized by its scalar memory, scalar update mechanism, and incorporation of memory mixing; and (ii) the mLSTM (detailed in Section~\ref{sec:mLSTM}), distinguished by a matrix-based memory and an update rule leveraging covariance (outer product), allowing complete parallelization. Both sLSTM and mLSTM improve upon the LSTM through the application of exponential gating. To facilitate parallel computation, memory mixing, specifically the recurrent hidden-to-hidden connections, is excluded from the mLSTM design. Both variants also allow for the expansion to multiple memory cells; within this arrangement, sLSTM enables memory mixing across cells. Additionally, sLSTM may include several heads without mixing memory between heads but permits mixing across cells within a single head, thus establishing a novel approach to memory mixing by combining exponential gating with the head structure. For mLSTM, the implementation of multiple heads is functionally identical to having multiple memory cells.

When these novel LSTM types are embedded within residual block modules, the resulting structures are referred to as xLSTM blocks (see Section~\ref{sec:xLSTMblock}). An xLSTM architecture is then realized by sequentially connecting such blocks in a residual fashion, as described in Section~\ref{sec:xLSTMarchitecure}. The components and overall organization of the xLSTM architecture are illustrated in Figure~\ref{fig:xlstm_sketch}.

\subsection{Review of the Long Short-Term Memory}
\label{sec:orgLSTM}

The original conception of LSTM~\citep{Hochreiter:91,Hochreiter:97nips,Hochreiter:97} presented the scalar memory cell as its core computational and storage element, designed to address the issue of vanishing gradients~\citep{Hochreiter:91,Hochreiter:00book} via the mechanism known as the constant error carousel (enabled by cell state updates). This memory cell operates with three primary gates: input, output, and forget gates. Notably, the forget gate was introduced by \citet{Gers:00}. The procedures for updating the LSTM memory cell at time step $t$ are given by:

\begin{align}
c_t \ &= \  f_t \ c_{t-1} \ + \ i_t \ z_t &  & &\text{cell state} \\
h_t  \ &= \ o_t \ \tilde{h}_t \ , & \tilde{h}_t \ &= \ \psi \left( c_t\right)
  &\text{hidden state} \\
z_t \ &= \ \varphi \left( \tilde{z}_t \right) \ , 
  &\tilde{z}_t \ &=  \ \mathbf{w}^\top_{z} \ \mathbf{x}_t \ + \
  r_{z}  h_{t-1} \ + \  b_{z} \ \
  &\text{cell input} \\
i_t \ &= \ \sigma \left( \tilde{i}_t  \right) \ , 
  &\tilde{i}_t \ &= \ \mathbf{w}^\top_{i} \ \mathbf{x}_t \ + \
  r_{i}  \ h_{t-1} \ + \  b_{i} \ \
  &\text{input gate} \\
f_t \ &= \ \sigma \left(  \tilde{f}_t \right) \ , 
  &\tilde{f}_t \ &= \ \mathbf{w}^\top_{f} \ \mathbf{x}_t  \ + \
  r_{f}  \ h_{t-1} \ + \  b_{f} \ \
  &\text{forget gate} \\
o_t \ &= \ \sigma \left( \tilde{o}_t \right) \ , 
  &\tilde{o}_t  \ &= \ \mathbf{w}^\top_{o} \ \mathbf{x}_t \ + \
  r_{o}  \ h_{t-1} \ + \  b_{o} \ \
  &\text{output gate} 
\end{align}

The input weight vectors $\mathbf{w}_{z}$, $\mathbf{w}_{i}$, $\mathbf{w}_{f}$, and $\mathbf{w}_{o}$ represent connections from the input vector $\mathbf{x}_t$ to the cell input, input gate, forget gate, and output gate, respectively. The parameters $r_{z}$, $r_{i}$, $r_{f}$, and $r_{o}$ denote the recurrent weights linking the previous hidden state $h_{t-1}$ to the cell input and each of the gates. Bias parameters $b_{z}$, $b_{i}$, $b_{f}$, and $b_{o}$ are associated with these respective components. The activation functions $\varphi$ and $\psi$ (commonly instantiated as $\tanh$) govern the transformations for the cell input and hidden state. In particular, $\psi$ serves to bound the cell state, which would not otherwise be limited. Sigmoid functions are employed for all gate activations, specifically, $\sigma \left( x \right)= 1/(1+\exp(-x))$. Later developments introduced the combination of several scalar memory cells into a vector $\mathbf{c}_t \in \mathbb{R}^d$, enabling the application of recurrent weight matrices $\mathbf{R} \in \mathbb{R}^{d \times d}$ to facilitate the interconnection of outputs between memory cells~\citep{Greff:15}. For further elaboration, see Appendix~\ref{sec:apporgLSTM}. Empirical ablation analyses have demonstrated the necessity of retaining all components within the memory cell architecture~\citep{Greff:15}.

\subsection{sLSTM}
\label{sec:sLSTM}

To enable LSTMs to adjust their storage choices, we incorporate exponential gates (highlighted in red), alongside mechanisms for normalization and stabilization. Specifically, the input and forget gates are allowed to utilize exponential activation functions. For normalization purposes, a separate normalizer state is added, which accumulates the sum of the input gate multiplied by all subsequent forget gates.

The scalar sLSTM forward pass is:

\begin{align}
\label{eq:slstmforward}
c_t \ &= \  f_t \ c_{t-1} \ + \ i_t \ z_t & & 
 &\text{cell state} \\
n_t \ &= \  f_t \ n_{t-1} \ + \ i_t  & & 
 &\text{normalizer state} \\
h_t \ &= \ o_t \ \tilde{h}_t \ ,
  &\tilde{h}_t \ &= \ c_t / n_t
&\text{hidden state} \\
z_t \ &= \ \varphi \left( \tilde{z}_t \right) \ , 
  &\tilde{z}_t \ &=  \ \mathbf{w}^\top_{z} \ \mathbf{x}_t \ + \
  r_{z}  h_{t-1} \ + \  b_{z}
  &\text{cell input} \\
i_t \ &= \ \!\exp \left( \tilde{i}_t  \right) \ , 
  &\tilde{i}_t \ &= \ \mathbf{w}^\top_{i} \ \mathbf{x}_t \ + \
  r_{i}  \ h_{t-1} \ + \  b_{i}
  &\text{input gate} \\
f_t \ &= \ \sigma \left(  \tilde{f}_t \right) \ \text{OR} \
\!\exp \left(  \tilde{f}_t \right) \ ,
  &\tilde{f}_t \ &= \ \mathbf{w}^\top_{f} \ \mathbf{x}_t  \ + \
  r_{f}  \ h_{t-1} \ + \  b_{f}
  &\text{forget gate} \\
o_t \ &= \ \sigma \left( \tilde{o}_t \right) \ , 
  &\tilde{o}_t  \ &= \ \mathbf{w}^\top_{o} \ \mathbf{x}_t \ + \
  r_{o}  \ h_{t-1} \ + \  b_{o}
  &\text{output gate} 
\end{align}

We transfer the original LSTM gating techniques, i.e., input- and/or hidden-dependent gating plus bias term, to the new architectures. Exponential activation functions can lead to large values that cause overflows. Therefore, we stabilize gates with an additional state \mbox{$m_t$~\citep{Milakov:18arxiv}}:

\begin{align}
\label{eq:slstmstabil}
m_t \ &= \ \max \left( \log ( f_t ) + m_{t-1} , \log ( i_t ) \right) &\text{stabilizer state} \\
i'_t \ &= \ \!\exp \left( \log \left ( i_t \right) - m_t \right) \ = \ \!\exp \left( \tilde{i}_t - m_t \right) \ \, 
  &\text{stabil. input gate} \\
f'_t \ &= \ \!\exp \left( \log \left( f_t \right) + m_{t-1} - m_t \right) \ \, 
  &\text{stabil. forget gate}
\end{align}

As demonstrated in Appendix~\ref{sec:appsLSTM}, substituting $f_t$ with $f'_t$ and $i_t$ with $i'_t$ during the forward propagation leaves both the network's overall output and the gradients of the loss function with respect to the model parameters unaffected.

\paragraph{New Memory Mixing.} Similar to the original LSTM, sLSTM is capable of maintaining several memory cells (refer to Appendix~\ref{sec:appsLSTM} for further details). The existence of multiple memory cells facilitates memory mixing through recurrent pathways $\mathbf{R}_{\mathbf{z}}$, $\mathbf{R}_{\mathbf{i}}$, $\mathbf{R}_{\mathbf{f}}$, $\mathbf{R}_{\mathbf{o}}$, which connect the hidden state vector $\mathbf{h}$ to the inputs of the memory cell $\mathbf{z}$ and the respective gates $\mathbf{i}$, $\mathbf{f}$, and $\mathbf{o}$. A distinguishing innovation in this approach to memory mixing lies in the application of exponential gating. With sLSTM, multiple heads can be incorporated, whereby memory mixing occurs within each head but does not extend between different heads. The combination of exponential gating and the use of multiple heads in sLSTM thus introduces a novel mechanism for mixing memory representations.

\subsection{mLSTM}
\label{sec:mLSTM}

In order to improve the storage capability of LSTMs, we reformulate the traditional memory cell, originally represented by a scalar $c \in \mathbb{R}$, as a matrix $\mathbf{C} \in \mathbb{R}^{d \times d}$. This adjustment means that memory retrieval is accomplished through matrix multiplication. At any time step $t$, we aim to memorize a vector pair—a key $\mathbf{k}_t \in \mathbb{R}^d$ and a value $\mathbf{v}_t \in \mathbb{R}^d$—adopting the terminology from the Transformer architecture. Subsequently, at a future time $t + \tau$, the corresponding value $\mathbf{v}_t$ should be accessed using a query vector $\mathbf{q}_{t+\tau} \in \mathbb{R}^d$. This formulation aligns with the framework established by Bidirectional Associative Memories (BAMs) \citep{Kohonen:72,Anderson:72,Nakano:72,Anderson:77}.

The covariance update rule~\citep{Sejnowski:77,Dayan:91} for storing a key-value pair is

\begin{align}
\mathbf{C}_t \ &= \ \mathbf{C}_{t-1} \ + \ \mathbf{v}_{t} \ \mathbf{k}_t^\top \ .
\end{align}

We posit that a layer normalization is performed prior to the projection of inputs onto keys and values, resulting in these vectors possessing a mean of zero. The covariance update method is provably optimal~\citep{Dayan:91} for maximizing the distinguishability between retrieved binary vectors, an objective synonymous with maximizing the signal-to-noise ratio. Enhanced separability can be attained if retrieval is confined to pairwise associations, although this entails accepting the quadratic computational cost characteristic of attention mechanisms~\citep{Krotov:16,Krotov:17,Ramsauer:21}. This covariance update is mathematically analogous to Fast Weight Programmers \citep{Schmidhuber:92ncfastweights,Schlag:21}, which have subsequently been extended to include a fixed decay coefficient applied to $\mathbf{C}_{t-1}$ and a static learning rate scaling 
$\mathbf{v}_{t} \mathbf{k}_t ^\top$~\citep{Ba:16}. Adopting this viewpoint, we embed the covariance update strategy within the LSTM architecture by mapping the forget gate to the decay coefficient, the input gate to the learning rate, and assigning scaling of the retrieved vector to the output gate.

For this matrix memory, the normalizer state is the weighted sum of key vectors, where each key vector is weighted by the input gate and all future forget gates. Again, the normalizer state keeps record of the strength of the gates. Since the dot product between query and normalizer state can be close to zero, we use the absolute value of this dot product and lower bound it by a threshold (typically 1.0) as done previously~\citep{Sun:23arxiv}. The mLSTM forward pass is:

\begin{align}
\label{eq:mlstm_recurrent_begin}
\mathbf{C}_t \ &= \  f_t \ \mathbf{C}_{t-1} \ + \ 
  i_t \ \mathbf{v}_t \ \mathbf{k}_t^\top & &  &\text{cell state} \\
\mathbf{n}_t \ &= \  f_t \ \mathbf{n}_{t-1} \ + \ 
  i_t \ \mathbf{k}_t & &  &\text{normalizer state} \\
\mathbf{h}_t  \ &= \ \mathbf{o}_t \ \odot \ \tilde{\mathbf{h}}_t
\ , \qquad \qquad 
\tilde{\mathbf{h}}_t \ = \   \mathbf{C}_t \mathbf{q}_t \ / \ 
  \max \left\{ |\mathbf{n}_t^\top \mathbf{q}_t|, 1 \right\} 
   &&&\text{hidden state} \\
\mathbf{q}_t \ &= \ \mathbf{W}_q \ \mathbf{x}_t \ + \ \mathbf{b}_q  & &  &\text{query input} \\
\mathbf{k}_t \ &= \ \frac{1}{\sqrt{d}} \mathbf{W}_k \ \mathbf{x}_t \ + \ \mathbf{b}_k  & &  &\text{key input} \\
\mathbf{v}_t \ &= \ \mathbf{W}_v \ \mathbf{x}_t \ + \ \mathbf{b}_v  & &  &\text{value input} \\
i_t \ &= \ \!\exp \!  \! \left( \tilde{i}_t  \right) \ , 
 \qquad \qquad \ \ \ \ \,
  \tilde{i}_t \ = \ \mathbf{w}^\top_{i} \ \mathbf{x}_t \ + \  b_{i} \ \ \ 
  &&&\text{input gate} \\
f_t \ &= \ \sigma \!  \left(  \tilde{f}_t \right) \ \text{OR} \ \!\exp \!  \! \left(  \tilde{f}_t \right) , \ \ \: \!
\tilde{f}_t \ = \ \mathbf{w}^\top_{f} \ \mathbf{x}_t  \ + \
  b_{f}
  &&& \text{forget gate} \\
\label{eq:mlstm_recurrent_end}
\mathbf{o}_t \ &= \ \sigma \left( \tilde{\mathbf{o}}_t \right) \ , \qquad \qquad \quad \ \ \ \,
  \tilde{\mathbf{o}}_t  \ = \ \mathbf{W}_{\mathbf{o}} \ \mathbf{x}_t \ + \
  \mathbf{b}_{\mathbf{o}}
  &&&\text{output gate} 
\end{align}

The mLSTM architecture allows for several memory cells, similar to the standard LSTM model. In the case of mLSTM, having multiple heads functions in the same way as having multiple cells due to the absence of memory mixing. To address the potentially unstable behavior of mLSTM's exponential gating, we apply the same stabilization methods utilized for sLSTM (refer to Equation~\ref{eq:slstmstabil}). 
Owing to the lack of memory mixing in mLSTM, its recurrent computation can be alternatively expressed in a parallel form. Additional information can be found in Appendix~\ref{sec:appmLSTM}.

\subsection{xLSTM Architecture}
\label{sec:xLSTMarch}

\paragraph{xLSTM Blocks.}
\label{sec:xLSTMblock}
An xLSTM block is designed to encapsulate past information in a nonlinear, high-dimensional representation, with the goal of improving the ability to distinguish between various contextual sequences. The clear separation of sequence histories is essential for accurately forecasting upcoming elements, such as subsequent tokens. This approach draws on Cover’s Theorem~\citep{Cover:65}, which posits that patterns transformed into higher-dimensional spaces via nonlinear mapping are, with higher probability, linearly separable compared to their original lower-dimensional counterparts. We explore two types of residual block designs: (i) The post up-projection residual block, commonly found in Transformers, which first applies nonlinear summarization in the input space before performing a linear up-projection to a higher-dimensional space, introducing a nonlinear activation, and finally mapping back linearly to the original dimension; refer to the leftmost panel of Figure~\ref{fig:Backbones} as well as the third column in Figure~\ref{fig:xlstm_sketch}. A more comprehensive illustration can be found in Figure~\ref{fig:appsLSTM_detailed} within the appendix. (ii) The pre up-projection residual block, similar in construction to State Space Models, where the process begins with a linear mapping into a higher-dimensional space,

compresses historical information through non-linear transformations within a high-dimensional latent space, followed by a linear projection back to the original feature space. In configurations where the xLSTM block utilizes an sLSTM, the up-projection occurs after the non-linear transformation. Conversely, for xLSTM blocks with an mLSTM, the up-projection precedes the non-linear operation to better accommodate the expanded memory capabilities afforded by the higher-dimensional space. Additional clarification can be found in the leftmost panel of Figure~\ref{fig:Backbones}, the third column of Figure~\ref{fig:xlstm_sketch}, or Figure~\ref{fig:appsLSTM_detailed} in the appendix.

\paragraph{xLSTM Architecture.}
\label{sec:xLSTMarchitecure}
The architecture of xLSTM is formed by residually stacking fundamental units, as outlined by~\citep{Srivastava:15,He:16}. Our implementation follows the prevalent pre-LayerNorm~\citep{Ba:16b} residual configuration, similar to those used in modern Large Language Models. Refer to the rightmost column of Figure~\ref{fig:xlstm_sketch} for architectural visualization.

\subsection{Memory and Speed Considerations}

Unlike Transformers, xLSTM networks exhibit linear computational requirements and maintain fixed memory usage relative to sequence length. Due to its ability to compress memory, xLSTM is highly appropriate for industrial use cases and deployment on edge devices.

In mLSTM, memory does not require learnable parameters, yet computations remain intensive because of the $d \times d$ matrix-based memory and the necessity for $d \times d$ updates. This design involves a trade-off, exchanging expanded memory capability for heightened computational load. Importantly, these computational demands are mitigated by parallel execution on GPUs, which limits their impact on actual processing time.

Although mLSTM supports parallelism in a manner comparable to FlashAttention~\citep{Dao:22,Dao:23} and GLA~\citep{Yang:23arxiv}, sLSTM is inherently non-parallelizable due to the mixing that arises from hidden-to-hidden connections in memory. Nonetheless, we introduced an efficient CUDA implementation for sLSTM with optimizations down to the GPU register level; as a result, sLSTM runs at a speed typically less than double that of mLSTM.

\section{Related Work}
\label{sec:related}

\paragraph{Linear Attention.} To address the Transformer’s quadratic time complexity with respect to sequence length, multiple approaches have been developed that achieve attention with linear complexity in context length. The Synthesizer, for example, generates synthetic attention scores independently of any token-to-token dependencies~\citep{Tay:20arxiv}. The Linformer approximates self-attention using a projection into a lower-rank space, making computation more efficient by reducing dimensionality~\citep{Wang:20arxiv}. Linear Transformer reconfigures the attention process into a linear operation~\citep{Katharopoulos:20}. Performer further optimizes attention by leveraging positive orthogonal random features to linearly estimate the softmax operation~\citep{Choromanski:21}. Alternatives such as the Structured Global Convolution (SGConv)~\citep{Li:22} and the Hyena Hierarchy~\citep{Poli:23} substitute attention entirely, employing long-sequence convolution mechanisms for fast global context aggregation.

\paragraph{State Space Models.} State Space Models (SSMs) have gained considerable attention recently due to their linear computational complexity with respect to context length and their competitive results when measured against Transformers. Notable early work in this direction includes the Structured State Space sequence model (S4)~\citep{Gu:21}. Subsequent developments introduced variations such as the Diagonal State Space (DSS) model~\citep{Gupta:22}, Gated State Space (GSS) models~\citep{Mehta:22}, the S5 model~\citep{Smith:22}, Bidirectional Gated SSM (BiGS)~\citep{Wang:22}, H3 model~\citep{Fu:23}, and more recently, Mamba~\citep{Gu:24arxiv}.

\paragraph{Recurrent Neural Networks.} Recurrent Neural Networks (RNNs) have been proposed as alternatives to Transformer architectures and attention mechanisms, mainly because of their linear scaling with respect to context length. Models utilizing Deep Linear Recurrent Units (LRUs) have achieved noteworthy performance in language modeling tasks~\citep{Orvieto:23, De:24arxiv}, and similar advancements have been observed with the Hierarchically Gated Linear RNN (HGRN)~\citep{Qin:23} as well as HGRN2~\citep{Qin:24arxiv}. Among RNN-based strategies for large language modeling, RWKV~\citep{Peng:23arxivshort,Peng:24arxivshort} has garnered attention, demonstrating results on par with leading Transformer models.

\paragraph{Gating.}
A central concept introduced by LSTM is the gating mechanism, a notion that has subsequently been adopted and re-envisioned in various recent architectures. Gating mechanisms feature prominently in models such as HGRN~\citep{Qin:23}, HGRN2~\citep{Qin:24arxiv}, Gated Linear Attention (GLA)~\citep{Yang:23arxiv}, Gated State Space (GSS) models~\citep{Mehta:22}, Bidirectional Gated SSM (BiGS)~\citep{Wang:22}, Moving Average Equipped Gated Attention (MEGA)~\citep{Ma:22}, RWKV~\citep{Peng:23arxivshort}, as well as in Mamba~\citep{Gu:24arxiv}.

\paragraph{Covariance Update Rule.} In order to extend memory capacity, the mLSTM cell was augmented with a matrix memory employing a covariance update rule. Comparable update techniques have also been utilized in Fast Weight Programmers \citep{Schmidhuber:92ncfastweights,Schlag:21}, RWKV-5 and RWKV-6~\citep{Peng:24arxivshort}, Retention~\citep{Sun:23arxiv}, Linear Transformer~\citep{Katharopoulos:20}, and HGRN2~\citep{Qin:24arxiv}.

\paragraph{Most Related.} The models most akin to xLSTM from a conceptual standpoint include Retention~\citep{Sun:23arxiv}, RWKV~\citep{Peng:23arxivshort,Peng:24arxivshort}, and HGRN2~\citep{Qin:24arxiv}, all of which incorporate either matrix memory, gating mechanisms, or both. Nevertheless, unlike the novel sLSTM, these models lack support for memory mixing. The ability to mix memory is instrumental for addressing state tracking challenges, rendering LSTMs more expressive than both State Space Models (SSMs) and Transformers~\citep{Merrill:24,Deletang:23}. Tasks such as code evaluation or maintaining consistency of entities in lengthy narratives fundamentally require efficient state tracking.

\paragraph{Residually Stacking Architectures.} Consistent with most modern, large-scale deep learning models, xLSTM architectures are designed by stacking their constituent blocks with residual connections~\citep{Srivastava:15,He:16}.

The use of this stacking strategy, combined with residual pathways, has been crucial in facilitating the training of deep convolutional networks~\citep{He:16} as well as the Transformer architecture~\citep{Vaswani:17}. Transformers, in particular, serve as the foundation for Large Language Models (LLMs) including prominent examples such as GPT-3~\citep{Brown:20short}, ChatGPT~\citep{Schulman:22short}, GPT-4~\citep{Achiam:23gpt4short}, Megatron-LM~\citep{Shoeybi:19}, Gopher~\citep{Rae:21short}, ERNIE 3.0 Titan~\citep{Wang:21arxivshort}, GLaM ~\citep{Du:21glamshort}, the Chinese M6~\citep{Lin:21}, multilingual AlexaTM 20B~\citep{Soltan:22}, OPT~\citep{Zhang:22opt}, Chinchilla~\citep{Hoffmann:22short}, BLOOM~\citep{Scao:22arxivshort}, GLM-130B~\citep{Zeng:22arxivshort}, LaMDA~\citep{Thoppilan:22short}, PaLM~\citep{Chowdhery:22short}, Llama~\citep{Touvron:23arxiv}, and Gemini~\citep{GeminiTeam:23arxiv,Reid:24arxiv}.

\section{Experiments}
\label{sec:Exp}

In this section, we present a series of empirical evaluations of xLSTM, comparing its performance to that of established approaches, with an emphasis on its application to language modeling. Section~\ref{sec:ExpSyntethic} is dedicated to probing the unique strengths of xLSTM on controlled synthetic benchmarks. 
Section~\ref{sec:ExpComparison} details a head-to-head comparison of various leading language modeling techniques using validation perplexity, following training on 15B tokens from the SlimPajama dataset~\citep{Soboleva:23}. Within this same experimental framework, we also perform ablations to dissect the internal mechanisms of xLSTM. Additionally, we analyze the scaling laws exhibited by these methods, in line with the analyses of \citet{Kaplan:20} and \citet{Brown:20short}.

Section~\ref{sec:ExpLanguage} extends this comparison to a larger scale. Here, we revisit xLSTM and the top-performing methods from Section~\ref{sec:ExpComparison}, this time training them on 300B tokens from SlimPajama~\citep{Soboleva:23}. The experimental analysis consists of: (1) measuring each method's generalization to longer context lengths, (2) benchmarking them using both validation perplexity and downstream task performance~\citep{Sutawika:24}, (3) assessing model performance across 571 textual domains from the PALOMA language benchmark~\citep{Magnusson:23arxivshort}, and (4) re-examining the scaling characteristics, now in the context of substantially increased (20x) training data.

\label{sec:xlstm_blocks_notation}
Throughout our experiments, we denote xLSTM[$a$:$b$] to specify the proportion $a/b$ between xLSTM blocks implemented using mLSTM and those implemented using sLSTM. For instance, xLSTM[7:1] designates a configuration in which, out of eight blocks, seven utilize mLSTM while one employs sLSTM. In setups containing 48 blocks in total, this configuration corresponds to 42 mLSTM-based blocks and 6 sLSTM-based blocks. Additionally, across all experimental setups, we apply pre up-projection to mLSTM blocks and post up-projection to sLSTM blocks.

\subsection{Synthetic Tasks and Long Range Arena}
\label{sec:ExpSyntethic}
We begin by evaluating the performance of xLSTM's novel exponential gating mechanism with memory mixing on formal language tasks~\citep{Deletang:23}. Subsequently, we examine the capabilities of xLSTM's advanced matrix memory through the Multi-Query Associative Recall task~\citep{Arora:23arxiv}. Lastly, we measure xLSTM's ability to handle long sequence processing by benchmarking it on the Long Range Arena~\citep{Tay:21}.

\paragraph{Evaluation of xLSTM's Exponential Gating Combined with Memory Mixing.} We examine the effectiveness of xLSTM's exponential gating mechanism augmented with memory mixing, designed to tackle state tracking tasks as discussed in~\citet{Merrill:24, Merrill:23}. To facilitate the investigation, we adapt and expand the formal language experiments introduced by \citet{Deletang:23} to permit training across sequences of multiple lengths for the purpose of length generalization. Comprehensive details of all the specific tasks as well as extended experimental outcomes are provided in Appendix~\ref{sec:appExpSynthetic-formalLang}. In our comparison, xLSTM is evaluated against alternative architectures, including Transformers, State Space Models, and various forms of Recurrent Neural Networks. Performance is quantified by measuring accuracy exclusively on tokens pertinent to the target task, with values normalized from 0 (reflecting chance) to 1 (indicating flawless performance). We benchmark the following two-block architectures: xLSTM[0:1] (sLSTM only), xLSTM[1:0] (mLSTM only), xLSTM[1:1], alongside Llama, Mamba, RWKV, Retention, Hyena, LSTM, and a Transformer employing LSTM within its blocks (LSTM (Block)). Figure~\ref{fig:formal-main} presents the experimental results. Notably, architectures like Transformers or State Space Models that lack the capacity for memory mixing (and thus, state tracking) fail on tasks involving regular grammars, such as the parity task. This finding is consistent with previous reports that such models are intrinsically less expressive than RNNs~\citep{Merrill:24, Merrill:23, Deletang:23}.

\paragraph{Test of xLSTM's Memory Capacities on Associative Recall Tasks.} This study examines the memory capacity of xLSTM's newly proposed matrix memory using the Multi-Query Associative Recall task~\citep{Arora:23arxiv}. In this setting, each sequence consists of randomly sampled key-value pairs drawn from an extensive vocabulary, which the models must retain for future retrieval. To intensify the challenge relative to the original benchmark, we increase the maximum number of key-value pairs to 256 and extend the context window up to 2048, thereby creating a more demanding evaluation of model memory capacity. Our comparisons span 2-block variants of Llama, Mamba, RWKV-5, RWKV-6, as well as xLSTM[1:1] and xLSTM[1:0]. Performance is quantified by accuracy in correctly retrieving stored key-value pairs. As Transformers (such as Llama) possess memory capacity that scales exponentially with their coding dimension~\citep{Ramsauer:21}, they are considered the benchmark for this task. The main results are visualized in Figure~\ref{fig:mqar-main}.

Among the evaluated non-Transformer models, xLSTM[1:1] outperforms all others, including under constrained model sizes. Notably, the inclusion of the sLSTM block does not impair memory performance; instead, it contributes positively, particularly evident in scenarios involving 256 key-value pairs, the most challenging test. Further results, discussed in Appendix~\ref{sec:appExpSynthetic-mqar}, involve extrapolation experiments demonstrating that the augmented memory structure of xLSTM enables generalization to context lengths surpassing those used during model training.

\paragraph{Test of xLSTM's Long Context Capabilities on Long Range Arena.} To evaluate how xLSTM manages extensive sequences and broader contextual information, we benchmark various approaches using the Long Range Arena~\citep{Tay:21}. Across all evaluated tasks, xLSTM achieves robust and reliable results, indicating that its architectural design is particularly adept at addressing the challenges inherent to long context scenarios.

For more details, see Appendix~\ref{sec:appExpSynthetic-lra}.

\subsection{Method Comparison and Ablation Study}
\label{sec:ExpComparison}

In order to investigate the central research question—specifically, the performance of our novel LSTM variants when scaled for language modeling—we conduct experiments by training xLSTMs alongside Transformers, State Space Models, and additional techniques, all on 15B tokens drawn from the SlimPajama dataset under an identical auto-regressive framework. The resulting models are then evaluated on the validation set, with supplementary ablation analyses conducted for the xLSTM architectures.

\paragraph{Comparison of xLSTM Against Alternative Approaches.} To provide a benchmark, we train each model using a dataset comprising 15B tokens extracted from SlimPajama~\citep{Soboleva:23}, and assess the models based on their perplexity measured over the validation split. The range of techniques evaluated includes our proposed xLSTM, as well as GPT-3 (Transformer architecture)~\citep{Brown:20short}, Llama (Transformer)~\citep{Touvron:23arxiv}, H3 (State Space Model, or SSM)~\citep{Fu:23}, Mamba (SSM)~\citep{Gu:24arxiv}, several RNN-based variants—RWKV-4, RWKV-5, RWKV-6~\citep{Peng:23arxivshort,Peng:24arxivshort}—in addition to GLA (a linear Transformer model)~\citep{Yang:23arxiv}, HGRN and HGRN2 (RNNs)~\citep{Qin:23,Qin:24arxiv}, RetNet (linear Transformer)~\citep{Sun:23arxiv}, Hyena (linear Transformer)~\citep{Poli:23}, and both xLSTM[1:0] and xLSTM[7:1] configurations (refer to Section~\ref{sec:xlstm_blocks_notation}). Model training primarily used mixed precision; however, in the cases of RWKV-5, RWKV-6, GLA, and HGRN2, mixed-precision training was implemented without PyTorch's automated mixed precision (see Appendix Section~\ref{sec:appGenTrainProc} for more). 

The evaluated methods are grouped into three broad categories: (a) Transformers, (b) State Space Models (SSMs), and (c) a joint group comprising Recurrent Neural Networks (RNNs) and linear Transformers, where linear Transformers replace conventional attention with linear mechanisms. Each model configuration maintains rough parity with a GPT-3 setup featuring 350M parameters, consisting of an embedding dimension of 1024 and 24 residual blocks. Notably, GPT-3 is distinct in its use of shared parameters for token and output embeddings, accounting for a reduced parameter count. According to the results summarized in Table~\ref{tab:spaj15b_model_results}, xLSTM achieves superior validation perplexity compared to the other baselines. More thorough experimental comparisons can be found in Appendix~\ref{sec:appExpComparison}. The scaling trends depicted in Figure~\ref{fig:temp_sclaw15B} suggest that xLSTM is also likely to excel as model size increases.

\paragraph{Ablation Studies.}
As illustrated in Table~\ref{tab:spaj15b_model_results} and Figure~\ref{fig:temp_sclaw15B}, the xLSTM model performs remarkably well in language modeling tasks when trained on 15B tokens from the SlimPajama dataset. To systematically examine the contributions of xLSTM's components, we incrementally adapt a standard LSTM model into the xLSTM configuration. The process begins by placing LSTM layers within pre-LayerNorm residual frameworks. Next, we modify the architecture to include a post up-projection stage. Lastly, we incorporate exponential gating mechanisms and matrix memory to complete the transformation.

The results are shown in Table~\ref{tab:ablstudies} (top). 

Through ablation studies, we ascribe the notable gains in performance to the combined effects of exponential gating and the matrix memory. Given the pivotal role of gating within RNNs and State Space Models, we conduct further ablations across various gating approaches. The findings summarized in Table~\ref{tab:ablstudies} (bottom) demonstrate that allowing each gate to be trainable and input-dependent yields cumulative improvements. Supplementary ablation analyses focused on particular backbone modules are presented in Appendix~\ref{sec:appAblDetails}.

\subsection{xLSTM as Large Language Model}
\label{sec:ExpLanguage}

This work concludes with comprehensive experiments in large-scale language modeling to evaluate the feasibility of deploying xLSTM as a large language model (LLM). To this end, we expand the training regime to encompass 300B tokens derived from SlimPajama, consistent with the token counts utilized in Mamba~\citep{Gu:24arxiv} and Griffin~\citep{De:24arxiv}.

For benchmarking, xLSTM is compared against RWKV-4, Llama, and Mamba—each chosen as representatives of their respective model families as outlined in Section~\ref{sec:ExpComparison}. RWKV-4 is selected to represent RNN-based approaches since robust training precision settings for RWKV-5, RWKV-6, and HGRN2 became available only after 300B token training commenced (details provided in Appendix~\ref{sec:appGenTrainProc}).

We train multiple model scales (125M, 350M, 760M, 1.3B), evaluating all models for their capacity to extrapolate to longer sequence lengths and measuring their validation set performance. Their effectiveness on downstream tasks is examined, as is their capability in language modeling across 571 domains featured in the PALOMA benchmark. Finally, we probe their scaling law characteristics.

\paragraph{Sequence Length Extrapolation.}
\label{sec:spaj300B_ctx_extrapolation}
We begin by evaluating the ability of 1.3B-parameter versions of xLSTM, RWKV-4, Llama, and Mamba to handle sequences beyond their training context. Each model is trained with a context window of 2048 tokens, and is subsequently assessed on sequences extending up to 16384 tokens. The empirical findings are summarized in Figure~\ref{fig:spaj300B_seq_extrapolation}. Notably, xLSTM consistently sustains low perplexity across extended sequence lengths, distinguishing itself from the other approaches.

\paragraph{Validation Perplexity and Downstream Tasks.}
\label{sec:spaj_downstream_eval}
Next, for every model scale, we measure the capabilities of xLSTM, RWKV-4, Llama, and Mamba on the SlimPajama validation dataset in the context of next token prediction, as well as on a suite of downstream benchmarks designed to evaluate common sense reasoning skills.

The third column in Table~\ref{tab:lmeval} presents the validation set perplexities obtained by the various approaches. For all model sizes, xLSTM[1:0] and xLSTM[7:1] consistently achieve the lowest validation set perplexities, marking them as the leading models in this respect. The remaining columns in Table~\ref{tab:lmeval} show results for several downstream tasks. In nearly every task and model size evaluated, xLSTM outperforms competing methods; an exception is seen with the ARC task, where Mamba occasionally attains superior performance. Further information is available in Appendix~\ref{sec:appExpLanguage}.

\paragraph{Performance on PALOMA Language Tasks.} Third, we evaluate models of varying sizes—including xLSTM, RWKV-4, Llama, and Mamba—on their ability to predict the next token within the suite of PALOMA language tasks~\citep{Magnusson:23arxivshort}. Evaluation is conducted by reporting the perplexity of next-token predictions over 571 distinct text domains, encompassing sources such as nytimes.com and Reddit's r/depression community.

Table~\ref{tab:ppleval} presents the results, organizing perplexity figures by both broad language modeling categories (first seven columns) and more granular domain-specific benchmarks (last five columns). Among the xLSTM variants, xLSTM[1:0] demonstrates superior performance on these language modeling tasks when compared to xLSTM[7:1].

Across the 571 evaluated text domains, xLSTM[1:0] achieves lower perplexity relative to Mamba on 568 occasions (99.5\%), to Llama on 486 domains (85.1\%), and to RWKV-4 on 570 domains (99.8\%). Expanded results and additional analysis can be found in Appendix~\ref{sec:appExpLanguage}.

\paragraph{Scaling Laws.} Fourth, we analyze the observed power-law scaling behavior, which enables projection of model performance to increased sizes~\citep{Kaplan:20,Brown:20short}. The scaling trends are depicted in Figure~\ref{fig:temp_sclaw300B}. Each model demonstrates a comparable pattern of scaling, though their baselines differ. Among these, RWKV-4 exhibits the lowest performance, while Llama and Mamba follow in ascending order. xLSTM consistently surpasses Mamba, and the gap in their performance mirrors the difference observed between Mamba and Llama. These scaling results suggest that as model size grows, xLSTM is likely to maintain or even improve its advantage relative to both Transformers and State-Space models.

\textbf{Generation Times and Maximal Throughput.} We present measurements of the time required for text generation in Figure~\ref{fig:inference_time_generation}, and report the maximum throughput in Figure~\ref{fig:inference_time_decoding_throughput} (left), focusing on the xLSTM models at the 1.3B parameter scale. Our analysis includes comparisons with equivalent Mamba, Llama, and RWKV models available via HuggingFace, utilizing a static key-value cache for Llama. During these experiments, compilation of the full cache for the Transformer Model and usage of \texttt{torch.compile} for compiling Mamba were not operational. For the text generation assessment, each model was benchmarked with a batch size of 1 and a pre-fill value of 16, which optimally benefits the Transformer.

The results in Figure~\ref{fig:inference_time_generation} demonstrate that xLSTM, along with other recurrent architectures such as Mamba and RWKV-4, exhibit linear time complexity in contrast to the quadratic scaling observed for Llama. Throughput measurements for decoding include various batch sizes and prefill parameters for the Llama baseline. As shown in Figure~\ref{fig:inference_time_decoding_throughput} (right), the constant memory footprint of xLSTM permits utilization of considerably larger batch sizes than Llama, leading to superior throughput.

\section{Limitations}

(i) Unlike mLSTM, the structure of sLSTM involves memory mixing that inherently prevents efficient parallel processing and, as a result, does not support rapid parallel execution. Despite this, we have engineered an optimized CUDA kernel for sLSTM, yielding performance that is currently less than twice as slow as our parallel version of mLSTM.

(ii) Our mLSTM CUDA kernels have not undergone thorough optimization, making our current solution roughly four times slower than implementations like FlashAttention or the scan approach used in Mamba. Additional speedups could potentially be realized by developing faster CUDA kernels similar to those used in FlashAttention.

(iii) The computational requirements for handling the matrix memory in mLSTM are considerable, as all $d \times d$ matrices need to be computed. Nevertheless, both the update and retrieval of memory are parameter-free and amenable to standard matrix-based parallelization, so the increased complexity of the memory module translates into only a minimal increase in actual computation time.

(iv) Careful consideration must be given to the initialization of the forget gates to ensure effective operation.

(v) Because the memory within the mLSTM is matrix-based and does not depend on the length of the input sequence, increasing the sequence length risks saturating memory resources, especially with longer contexts. However, this has not posed a limitation for contexts as long as 16k tokens (refer to Section~\ref{sec:spaj300B_ctx_extrapolation}).

(vi) Owing to the high computational demands of large-scale language modeling experiments, we have not extensively tuned the model architecture or hyperparameters—particularly for more sizable xLSTM variants. We expect that systematic optimization is essential for the xLSTM family to fully realize its capabilities.

\section{Conclusion}
We have provided a partial response to our central question: To what extent does language modeling improve when scaling LSTM architectures to billions of parameters? Our findings suggest that such scaling allows LSTMs to achieve performance on par with leading approaches such as Transformers and State Space Models. We introduced the xLSTM, augmenting LSTM with exponential gating featuring memory mixing, as well as an innovative memory architecture. In empirical evaluations of language modeling, xLSTM models show competitive or superior performance compared to cutting-edge frameworks like Transformers and State Space Models. The observed scaling behaviors imply that enlarging xLSTM models may position them as formidable alternatives to existing Large Language Models predominantly based on Transformer designs. Moreover, xLSTM holds significant promise for advancing domains including Reinforcement Learning, Time Series Prediction, and the modeling of physical systems.

\end{document}